% Source: Wolf (2012), Chapter 8, lines 1221--1224 and 1247--1254.

\begin{definition}[Jordan condition number]
    \label{def:wolf_jordan_condition_number}
    \lean{Matrix.IsJordanMatrix, Matrix.jordanConditionFactors, Matrix.jordanConditionNumber}
    \leanok
    A normalized Jordan matrix is a finite direct sum of blocks
    $J_{d_k}(\lambda_k)$ with $d_k>0$, up to a permutation of coordinates.
    For a complex square matrix $X$, let
    \begin{align}
      \mathcal C_X&=\{\lVert P\rVert_\infty\lVert Q\rVert_\infty:
        PQ=QP=I,\ X=PJQ,\ J\text{ is a normalized Jordan matrix}\},\\
      \kappa_X&=\inf\mathcal C_X.
    \end{align}
    This is the definition in \cite[Chapter~8]{Wolf2012Quantum}.
\end{definition}

\begin{theorem}[Admissible forms and the condition infimum]
    \label{thm:wolf_jordan_condition_infimum}
    \lean{Matrix.isJordanMatrix_diagonal, Matrix.isJordanMatrix_jordanBlockDirectSum,
      Matrix.bddBelow_jordanConditionFactors, Matrix.jordanConditionNumber_le,
      Matrix.one_le_jordanConditionNumber, Matrix.exists_jordanConditionFactor_lt,
      Matrix.jordanConditionNumber_of_isEmpty, Matrix.one_le_l2_opNorm_mul_of_mul_eq_one}
    \leanok
    Diagonal matrices and finite sums of positive-dimensional Jordan blocks
    are normalized Jordan matrices. The family $\mathcal C_X$ is bounded
    below, and $\kappa_X\le c$ for every $c\in\mathcal C_X$.
    If the underlying space has positive dimension and $\mathcal C_X$ is
    nonempty, then $1\le\kappa_X$. If $\mathcal C_X$ is nonempty, then
    for every $\varepsilon>0$ there is $c\in\mathcal C_X$ with
    $c<\kappa_X+\varepsilon$.
    On the zero-dimensional space, $\kappa_X=0$.
\end{theorem}
\begin{proof}
    \leanok
    \uses{def:wolf_jordan_condition_number}
    A diagonal matrix is the sum of its one-dimensional blocks.
    Submultiplicativity gives $1=\lVert PQ\rVert_\infty
    \le\lVert P\rVert_\infty\lVert Q\rVert_\infty$ on a nonzero space.
    The greatest-lower-bound property proves the inequalities and the
    $\varepsilon$ approximation. On the zero space all matrix norms vanish.
\end{proof}

\begin{theorem}[Uniform comparisons pass to the infimum]
    \label{thm:wolf_jordan_infimum_comparison}
    \lean{Matrix.jordanConditionNumber_bounds,
      Matrix.l2_opNorm_pow_jordanConditionNumber_bounds}
    \leanok
    Suppose $X$ acts on a nonzero finite-dimensional space and
    $\mathcal C_X\ne\varnothing$. Let $r,s\ge0$. If $r\le cs$ and $s\le cr$
    for every $c\in\mathcal C_X$, then $\kappa_X^{-1}r\le s\le\kappa_Xr$.
    In particular, if every admissible Jordan representative of $X$ has
    power norm $\lVert J^n\rVert_\infty=r$, then
    \begin{align}
      \kappa_X^{-1}r\le\lVert X^n\rVert_\infty\le\kappa_Xr.
    \end{align}
\end{theorem}
\begin{proof}
    \leanok
    \uses{thm:wolf_jordan_condition_infimum, lem:wolf_jordan_similarity_comparison}
    For $s>0$, $r/s$ is a lower bound for $\mathcal C_X$, hence
    $r/s\le\kappa_X$. For $r>0$, apply the same argument to $s/r$.
    If either is zero, use any member of the nonempty family.
    For each admissible $X=PJQ$ with $PQ=QP=I$, the similarity comparison gives
    \begin{align}
      (\lVert P\rVert_\infty\lVert Q\rVert_\infty)^{-1}
        \lVert J^n\rVert_\infty
      \le\lVert X^n\rVert_\infty
      \le\lVert P\rVert_\infty\lVert Q\rVert_\infty
        \lVert J^n\rVert_\infty.
    \end{align}
    Set $c=\lVert P\rVert_\infty\lVert Q\rVert_\infty>0$ and
    $s=\lVert X^n\rVert_\infty$. The common value
    $r=\lVert J^n\rVert_\infty$ therefore satisfies $r\le cs$ and $s\le cr$
    for every $c\in\mathcal C_X$, so the uniform comparison applies.
\end{proof}

\begin{theorem}[Scaling after a supplied Jordan similarity]
    \label{thm:wolf_supplied_jordan_similarity_scaling}
    \lean{Matrix.exists_pos_pow_similarity_jordanBlockDirectSum_bounds}
    \leanok
    Let $J=\bigoplus_k J_{d_k}(\lambda_k)$ be a finite sum with $d_k>0$,
    let $X=PJQ$ with $PQ=QP=I$, and set
    $c=\lVert P\rVert_\infty\lVert Q\rVert_\infty$.
    Choose a block of maximal modulus $0<\mu\le1$ and maximal dimension $d$
    among ties. There exists $C_2>0$ such that for every positive integer $n$,
    $\lVert X^n\rVert_\infty\le C_2\mu^n n^{d-1}$.
    For every $n\ge d-1$,
    \begin{align}
      c^{-1}(\mu(d-1))^{-(d-1)}\mu^n n^{d-1}
      \le\lVert X^n\rVert_\infty.
    \end{align}
    When $d=1$, the zeroth-power factor is one.
\end{theorem}
\begin{proof}
    \leanok
    \uses{lem:wolf_jordan_similarity_comparison,
      thm:wolf_finite_jordan_sum_polynomial_scaling}
    Multiply the finite-sum upper constant by $c$. Multiply its lower
    estimate by $c^{-1}$ and apply the similarity comparison.
\end{proof}
% The channel decomposition, identification of the subperipheral blocks and
% invariance of their power norm over all admissible forms remain separate
% obligations. This is not a completion of Wolf Theorem 8.23.


\begin{theorem}[Nilpotence after a supplied Jordan similarity]
    \label{thm:wolf_supplied_jordan_nilpotence}
    \lean{Matrix.pow_eq_zero_of_similarity_jordanBlockDirectSum_maximal_zero}
    \leanok
    Let $X=PJQ$, $PQ=QP=I$, with $J$ a nonempty finite sum of
    positive-dimensional Jordan blocks. If the maximal eigenvalue modulus
    is zero and the largest block dimension is $d$, then $X^n=0$ for $n\ge d$.
\end{theorem}
\begin{proof}
    \leanok
    \uses{thm:wolf_finite_jordan_sum_nilpotence, lem:wolf_jordan_similarity_comparison}
    The finite-sum nilpotence theorem gives $J^n=0$. The similarity upper
    bound gives $\lVert X^n\rVert_\infty\le0$, hence $X^n=0$.
\end{proof}
